\subsection{Absolute Value Functions} 
To better understand them, we organize functions into \term{families}.
\begin{definition}[Family of Functions]
A \term{family of functions} consists of a base function, composed with one or more linear functions.
\end{definition}
\begin{definition}[Absolute Value Function]
The \term{absolute value functions} are the family of functions that include $f(x)=\Abs{x}$.
\end{definition}
\begin{example}Let $f(x)=\Abs{x}$, $g(x)=2x-1$, and $h(x)=x+4$. Find the following.
\begin{lpicvsplit}{2}{3}
\regtextleft{1}{-1}{$(g\circ f\circ h)(x)={}$};
\regtextleft{\value{fcols}+2}{-1}{$(h\circ f\circ g)(x)={}$};
\end{lpicvsplit}
\end{example}
\begin{example}
For each function below, explain why it is \emph{not} an absolute value function.
\begin{itemize}
\item $f(x)=\Abs{\frac{1}{x}}$\makeblanknewf
\item $f(x)=\Abs{2x-1}+x$\makeblanknewf
\item $f(x)=x\Abs{4x+3}$\makeblanknewf
\item $f(x)=\sqrt{\Abs{2x}}$\makeblanknewf
\end{itemize}
\end{example}